\documentclass[11pt]{article}
\usepackage[margin=1in]{geometry}
\usepackage{xcolor}
\usepackage[breakable,listings]{tcolorbox}
\tcbuselibrary{listings}
\usepackage{hyperref}
\usepackage{enumitem}
\usepackage{listings}
\usepackage{amsmath}
\usepackage{booktabs}
\usepackage{array}

% Define colors
\definecolor{osaorange}{RGB}{220,115,38}
\definecolor{aiblue}{RGB}{30,144,255}

% Configure hyperref
\hypersetup{
    colorlinks=true,
    linkcolor=blue,
    urlcolor=blue
}

% Code environment
\newtcblisting{rcode}{
  colback=white,
  colframe=gray!50,
  boxrule=0.5pt,
  arc=0pt,
  left=3pt,
  right=3pt,
  top=3pt,
  bottom=3pt,
  width=\textwidth,
  breakable,
  listing only,
  listing options={
    basicstyle=\small\ttfamily,
    breaklines=true,
    breakatwhitespace=true,
    columns=flexible,
    keepspaces=true,
    showstringspaces=false
  }
}

\title{}
\author{}
\date{}

\begin{document}

% Header box
\begin{tcolorbox}[colback=osaorange,colframe=black,boxrule=2pt,arc=0pt,left=3pt,right=3pt,top=6pt,bottom=6pt,boxsep=2pt]
\centering
\textcolor{white}{\textbf{\Large H524 Introduction to Biostatistics}}\\
\textcolor{white}{\textbf{\Large Week 3: Probability Distributions and Sampling}}\\
\textcolor{white}{\textbf{\Large Homework Assignment}}\\
\textcolor{white}{\textbf{\Large Fall 2025}}
\end{tcolorbox}

\vspace{8pt}

\noindent\textbf{Due Date:} End of Week 4\\
\textbf{Total Points:} 100 (+ 5 bonus)\\
\textbf{Format:} Submit through Canvas

\section*{Instructions}

\begin{itemize}
\item Show all work for full credit
\item Round probabilities to 4 decimal places unless otherwise specified
\item Round percentages to 1 decimal place
\item Include R code where specified
\item You may use R to check your work
\end{itemize}

\section{Part 1: Binomial Distribution (25 points)}

\subsection*{Question 1: Binomial Conditions (5 points)}

A pharmaceutical company claims that their new medication is effective for 80\% of patients. Which statement BEST describes why treating 15 patients constitutes a binomial experiment?

\begin{enumerate}[label=\Alph*)]
\item Each patient has two outcomes (success/failure), probability stays constant at 0.80, trials are independent, and there are a fixed number (15) of trials
\item Each patient can have multiple outcomes, the probability changes with each patient, trials depend on previous results, and sample size is variable
\item The medication works most of the time, some patients get better, the trials happen sequentially, and n=15
\item Each patient is tested once, the success rate is high, patients don't affect each other, and we count total successes
\end{enumerate}

\vspace{0.5cm}

\subsection*{Question 2: Binomial Probability - Exactly 12 Successes (3 points)}

For the medication scenario ($n=15$, $p=0.80$), what is the probability that exactly 12 patients respond successfully?

\begin{enumerate}[label=\Alph*)]
\item 0.2501
\item 0.8000
\item 0.1876
\item 0.3980
\end{enumerate}

\vspace{0.5cm}

\subsection*{Question 3: Binomial Probability - At Least 13 Successes (3 points)}

For the same scenario, what is the probability that at least 13 patients respond successfully?

\begin{enumerate}[label=\Alph*)]
\item 0.6482
\item 0.3980
\item 0.2309
\item 0.5000
\end{enumerate}

\vspace{0.5cm}

\subsection*{Question 4: R Code for Binomial (2 points)}

Which R code correctly calculates $\Pr(X \geq 13)$ for $n=15$, $p=0.80$?

\begin{enumerate}[label=\Alph*)]
\item \texttt{sum(dbinom(13:15, 15, 0.80))} or \texttt{pbinom(12, 15, 0.80, lower.tail=FALSE)}
\item \texttt{dbinom(13, 15, 0.80)}
\item \texttt{pbinom(13, 15, 0.80)}
\item \texttt{1 - dbinom(0:12, 15, 0.80)}
\end{enumerate}

\vspace{0.5cm}

\subsection*{Question 5: Binomial Expected Value (2 points)}

For the binomial scenario above ($n=15$, $p=0.80$), what is the expected number of successful responses?

\begin{enumerate}[label=\Alph*)]
\item 12
\item 15
\item 10
\item 9.6
\end{enumerate}

\vspace{0.5cm}

\subsection*{Question 6: Binomial Standard Deviation (2 points)}

What is the standard deviation for this binomial distribution?

\begin{enumerate}[label=\Alph*)]
\item 1.549
\item 2.400
\item 12.000
\item 0.800
\end{enumerate}

\vspace{0.5cm}

\subsection*{Question 7: Vaccine Efficacy - Lower Tail (3 points)}

A vaccine has 90\% efficacy. A clinic vaccinates 25 people. What is the probability that fewer than 20 people are protected?

\begin{enumerate}[label=\Alph*)]
\item 0.0334
\item 0.1607
\item 0.9018
\item 0.5000
\end{enumerate}

\vspace{0.5cm}

\subsection*{Question 8: Vaccine Efficacy - Range (3 points)}

For the same scenario, what is the probability that between 22 and 24 people (inclusive) are protected?

\begin{enumerate}[label=\Alph*)]
\item 0.6918
\item 0.2659
\item 0.1994
\item 0.3826
\end{enumerate}

\vspace{0.5cm}

\section{Part 2: Normal Distribution (25 points)}

\subsection*{Question 9: Z-Score Calculation (2 points)}

Adult heights are normally distributed with mean 170 cm and standard deviation 10 cm. What is the Z-score for a person who is 185 cm tall?

\begin{enumerate}[label=\Alph*)]
\item 1.5
\item 0.15
\item 2.0
\item -1.5
\end{enumerate}

\vspace{0.5cm}

\subsection*{Question 10: Upper Tail Probability (3 points)}

Using the same height distribution, what proportion of adults are taller than 185 cm?

\begin{enumerate}[label=\Alph*)]
\item 0.0668
\item 0.9332
\item 0.1587
\item 0.5000
\end{enumerate}

\vspace{0.5cm}

\subsection*{Question 11: Probability Between Two Values (3 points)}

What proportion of adults have heights between 165 cm and 180 cm?

\begin{enumerate}[label=\Alph*)]
\item 0.5328
\item 0.3085
\item 0.6915
\item 0.8413
\end{enumerate}

\vspace{0.5cm}

\subsection*{Question 12: 90th Percentile (3 points)}

Using the height distribution (mean 170 cm, SD 10 cm), what height represents the 90th percentile?

\begin{enumerate}[label=\Alph*)]
\item 182.8 cm
\item 190.0 cm
\item 180.0 cm
\item 175.6 cm
\end{enumerate}

\vspace{0.5cm}

\subsection*{Question 13: 25th Percentile (2 points)}

What height represents the 25th percentile?

\begin{enumerate}[label=\Alph*)]
\item 163.3 cm
\item 165.0 cm
\item 167.5 cm
\item 160.0 cm
\end{enumerate}

\vspace{0.5cm}

\subsection*{Question 14: Empirical Rule Range (2 points)}

Cholesterol levels are normally distributed with mean 200 mg/dL and SD 40 mg/dL. Using the empirical rule, within what range do approximately 95\% of cholesterol values fall?

\begin{enumerate}[label=\Alph*)]
\item 120 to 280 mg/dL
\item 160 to 240 mg/dL
\item 180 to 220 mg/dL
\item 100 to 300 mg/dL
\end{enumerate}

\vspace{0.5cm}

\subsection*{Question 15: Extreme Cholesterol Values (3 points)}

What proportion have cholesterol levels above 280 mg/dL?

\begin{enumerate}[label=\Alph*)]
\item 0.0228
\item 0.9772
\item 0.0500
\item 0.1587
\end{enumerate}

\vspace{0.5cm}

\subsection*{Question 16: Hypertension Proportion (2 points)}

Systolic blood pressure is normally distributed with mean 125 mmHg and SD 18 mmHg. What proportion have blood pressure above 140 mmHg (hypertension)?

\begin{enumerate}[label=\Alph*)]
\item 0.2024
\item 0.7976
\item 0.1587
\item 0.0668
\end{enumerate}

\vspace{0.5cm}

\subsection*{Question 17: Expected Number Hypertensive (2 points)}

If 1000 people are screened, approximately how many would be classified as hypertensive?

\begin{enumerate}[label=\Alph*)]
\item 202
\item 798
\item 159
\item 140
\end{enumerate}

\vspace{0.5cm}

\newpage

\section{Part 3: Sampling Distributions and CLT (25 points)}

\subsection*{Question 18: Sampling Distribution Mean (2 points)}

A population has mean $\mu = 100$ and standard deviation $\sigma = 20$. If you take a sample of size $n=25$, what is the mean of the sampling distribution of $\bar{X}$?

\begin{enumerate}[label=\Alph*)]
\item 100
\item 20
\item 4
\item 25
\end{enumerate}

\vspace{0.5cm}

\subsection*{Question 19: Standard Error Calculation (2 points)}

What is the standard error of $\bar{X}$ for $n=25$?

\begin{enumerate}[label=\Alph*)]
\item 4
\item 20
\item 100
\item 0.8
\end{enumerate}

\vspace{0.5cm}

\subsection*{Question 20: Effect of Sample Size on SE (3 points)}

If $n=100$ instead of $n=25$, what is the new standard error?

\begin{enumerate}[label=\Alph*)]
\item 2
\item 4
\item 10
\item 1
\end{enumerate}

\vspace{0.5cm}

\subsection*{Question 21: CLT Distribution (2 points)}

A population of glucose levels has mean 100 mg/dL and standard deviation 25 mg/dL. A sample of 64 individuals is randomly selected. According to the Central Limit Theorem, the sample mean $\bar{X}$ is distributed as:

\begin{enumerate}[label=\Alph*)]
\item $N(100, 3.125^2)$
\item $N(100, 25^2)$
\item $N(100, 0.391^2)$
\item $N(64, 25^2)$
\end{enumerate}

\vspace{0.5cm}

\subsection*{Question 22: Sample Mean Probability - Upper Tail (4 points)}

For the glucose scenario above, what is the probability that the sample mean exceeds 105 mg/dL?

\begin{enumerate}[label=\Alph*)]
\item 0.0548
\item 0.4207
\item 0.9452
\item 0.5000
\end{enumerate}

\vspace{0.5cm}

\subsection*{Question 23: Sample Mean Probability - Range (2 points)}

What is the probability that the sample mean glucose level is between 97 and 103 mg/dL?

\begin{enumerate}[label=\Alph*)]
\item 0.6629
\item 0.3830
\item 0.4515
\item 0.7881
\end{enumerate}

\vspace{0.5cm}

\subsection*{Question 24: SD vs SE Calculation (2 points)}

If the population standard deviation is 30 and sample size is 36, what is the standard error?

\begin{enumerate}[label=\Alph*)]
\item 5
\item 30
\item 6
\item 1.2
\end{enumerate}

\vspace{0.5cm}

\subsection*{Question 25: Effect of Sample Size on SE (2 points)}

As sample size increases, what happens to the standard error and why is this important?

\begin{enumerate}[label=\Alph*)]
\item SE decreases; this makes our estimates more precise and reliable for inference
\item SE increases; this makes our estimates more conservative
\item SE stays the same; sample size doesn't affect variability
\item SE decreases; this makes the population standard deviation smaller
\end{enumerate}

\vspace{0.5cm}

\subsection*{Question 26: SD vs SE Concept (1 point)}

What is the key difference between standard deviation (SD) and standard error (SE)?

\begin{enumerate}[label=\Alph*)]
\item SD measures variability of individual observations; SE measures variability of sample means
\item SD measures variability of sample means; SE measures variability of individual observations
\item SD is always larger than SE
\item SD and SE are the same when n=1
\end{enumerate}

\vspace{0.5cm}

\subsection*{Question 27: CLT Birth Weight Probability (3 points)}

Birth weights have mean 3400g and SD 500g. You sample 100 births. What is the probability that the sample mean birth weight is less than 3350g?

\begin{enumerate}[label=\Alph*)]
\item 0.1587
\item 0.4602
\item 0.8413
\item 0.0228
\end{enumerate}

\vspace{0.5cm}

\subsection*{Question 28: R Code for CLT (2 points)}

Which R code correctly calculates the probability in Question 27?

\begin{enumerate}[label=\Alph*)]
\item \texttt{pnorm(3350, 3400, 500/sqrt(100))} or \texttt{pnorm(3350, 3400, 50)}
\item \texttt{pnorm(3350, 3400, 500)}
\item \texttt{pbinom(3350, 100, 0.5)}
\item \texttt{pnorm(3350, 3400, 50, lower.tail=FALSE)}
\end{enumerate}

\vspace{0.5cm}

\newpage

\section{Part 4: Confidence Intervals (20 points)}

\subsection*{Question 13: Manual CI Calculation (10 points)}

A study measures pain reduction scores for 20 patients receiving a new treatment:

Sample mean: $\bar{x} = 6.2$ points\\
Sample standard deviation: $s = 2.4$ points\\
Sample size: $n = 20$

\begin{enumerate}[label=\alph*)]
\item Should you use a $z$ or $t$ distribution? Why? (2 points)
\item Find the appropriate critical value for a 95\% confidence interval. (2 points)
\item Calculate the standard error. (2 points)
\item Calculate the margin of error. (2 points)
\item Construct the 95\% confidence interval. (2 points)
\end{enumerate}

\vspace{1.5cm}

\subsection*{Question 29: CI Interpretation (2 points)}

Based on Question 13, which statement correctly interprets the 95\% confidence interval for mean pain reduction?

\begin{enumerate}[label=\Alph*)]
\item We are 95\% confident that the true mean pain reduction for all patients receiving this treatment is between 5.08 and 7.32 points
\item 95\% of patients will have pain reduction between 5.08 and 7.32 points
\item There is a 95\% probability that the sample mean is between 5.08 and 7.32 points
\item The population mean is definitely between 5.08 and 7.32 points with 95\% certainty
\end{enumerate}

\vspace{0.5cm}

\subsection*{Question 30: Meaning of 95\% Confidence (2 points)}

What does "95\% confident" mean in the context of confidence intervals?

\begin{enumerate}[label=\Alph*)]
\item If we repeated this study many times, 95\% of the confidence intervals we construct would contain the true population mean
\item 95\% of the data values fall within the confidence interval
\item There is a 95\% probability that our specific interval contains the population mean
\item We are 95\% sure our sample mean equals the population mean
\end{enumerate}

\vspace{0.5cm}

\subsection*{Question 31: Comparing Confidence Levels (1 point)}

If you calculated a 99\% confidence interval instead of 95\%, the interval would be:

\begin{enumerate}[label=\Alph*)]
\item Wider, because we need a larger range to be more confident we've captured the true mean
\item Narrower, because we are more confident about the estimate
\item The same width, confidence level doesn't affect interval width
\item Wider, but only if sample size increases
\end{enumerate}

\vspace{0.5cm}

\subsection*{Question 15: CI with R (5 points)}

A clinical trial measures cholesterol reduction for 30 patients:

\texttt{cholesterol\_reduction <- c(28, 32, 25, 30, 27, 35, 22, 29, 31, 26,}\\
\texttt{24, 33, 28, 30, 27, 29, 31, 25, 28, 32,}\\
\texttt{26, 30, 29, 27, 34, 28, 31, 25, 29, 30)}

\begin{enumerate}[label=\alph*)]
\item Write R code to calculate a 95\% confidence interval using \texttt{t.test()}. (2 points)
\item Report the confidence interval. (1 point)
\item Based on this interval, is there strong evidence that the mean reduction exceeds 25 mg/dL? Explain. (2 points)
\end{enumerate}

\vspace{1.5cm}

\section{Part 5: Integration and R Applications (5 points)}

\subsection*{Question 16: Comprehensive Problem (5 points)}

A pharmaceutical company tests a new drug. In trials with 80 patients:
\begin{itemize}
\item 72 patients respond successfully
\item Mean cholesterol reduction: $\bar{x} = 35$ mg/dL
\item Standard deviation: $s = 8$ mg/dL
\end{itemize}

\begin{enumerate}[label=\alph*)]
\item Calculate a 95\% confidence interval for the true success rate (proportion). Hint: For proportions, use $\hat{p} \pm 1.96\sqrt{\frac{\hat{p}(1-\hat{p})}{n}}$ (2 points)
\item Calculate a 95\% confidence interval for the true mean cholesterol reduction. (2 points)
\item Write complete R code to perform both calculations and create a visualization of the cholesterol reduction data. (1 point)
\end{enumerate}

\vspace{1.5cm}

\section*{Bonus: AI-Enhanced Learning (5 bonus points)}

\subsection*{Question 17: AI Verification and Learning (5 bonus points)}

Choose ONE problem from Part 3 (Sampling Distributions and CLT).

\begin{enumerate}[label=\alph*)]
\item Ask an AI tool (Claude, ChatGPT, etc.) to solve the problem. Include your exact prompt. (1 point)
\item Verify the AI's answer by working through it yourself step-by-step. Did the AI get it correct? If not, identify the error. (2 points)
\item Ask the AI to explain the Central Limit Theorem using a medical example different from those in class. Evaluate the quality of the explanation. (2 points)
\end{enumerate}

\vspace{2cm}

\newpage

\section*{Submission Requirements}

\subsection*{Required Files:}
\begin{enumerate}
\item \textbf{PDF Document} with all written work, calculations, and interpretations
\item \textbf{R Script File (.R)} with all R code properly commented
\item \textbf{AI Usage Documentation} (if applicable for bonus):
   \begin{itemize}
   \item Prompts used
   \item AI responses
   \item Your verification and evaluation
   \end{itemize}
\end{enumerate}

\subsection*{Formatting Guidelines:}
\begin{itemize}
\item Clearly label each question number
\item Show all calculation steps for full credit
\item For probability problems, write both the formula and the numerical answer
\item Round consistently per instructions (probabilities to 4 decimals, percentages to 1 decimal)
\item Include units where appropriate (mg/dL, mmHg, cm, etc.)
\end{itemize}

\section*{Grading Rubric}

\begin{table}[h]
\centering
\begin{tabular}{lcp{8cm}}
\toprule
\textbf{Category} & \textbf{Points} & \textbf{Criteria} \\
\midrule
Calculations & 40 & Correct formulas, accurate arithmetic \\
Methods & 30 & Appropriate approach, proper formulas \\
Interpretation & 20 & Clear explanations, correct context \\
R Code & 10 & Working code, proper syntax, comments \\
Bonus & +5 & Thoughtful AI integration and evaluation \\
\bottomrule
\end{tabular}
\end{table}

\section*{Common Mistakes to Avoid}

\begin{itemize}
\item \textbf{Binomial:} Forgetting to check conditions; using continuous formulas for discrete distributions
\item \textbf{Normal:} Confusing $\Pr(X > a)$ with $\Pr(X < a)$; forgetting to standardize
\item \textbf{Sampling Distributions:} Using population SD instead of SE; confusing individual vs. sample mean distributions
\item \textbf{Confidence Intervals:} Using $z$ when you should use $t$; misinterpreting what "95\% confident" means
\item \textbf{R Code:} Using \texttt{lower.tail=TRUE} when you need \texttt{lower.tail=FALSE}; forgetting to specify \texttt{mean} and \texttt{sd} parameters
\end{itemize}

\section*{Key Formulas}

\textbf{Binomial Distribution:}
$$\Pr(X = k) = \binom{n}{k} p^k (1-p)^{n-k}, \quad E(X) = np, \quad \text{SD}(X) = \sqrt{np(1-p)}$$

\textbf{Standardization:}
$$Z = \frac{X - \mu}{\sigma}$$

\textbf{Sampling Distribution:}
$$E(\bar{X}) = \mu, \quad \text{SE}(\bar{X}) = \frac{\sigma}{\sqrt{n}}$$

\textbf{Confidence Interval (t-distribution):}
$$\bar{x} \pm t_{\alpha/2, n-1} \times \frac{s}{\sqrt{n}}$$

\section*{Academic Integrity}

\begin{tcolorbox}[colback=aiblue!10,colframe=aiblue,title=AI Usage Guidelines]

\textbf{You are encouraged to:}
\begin{itemize}
\item[+] Use AI tools to check your work
\item[+] Ask AI to explain concepts you don't understand
\item[+] Verify your calculations with AI
\item[+] Use AI to help debug R code
\end{itemize}

\textbf{You must NOT:}
\begin{itemize}
\item[-] Submit AI-generated work without understanding it
\item[-] Copy solutions without showing your own reasoning
\item[-] Use AI as a substitute for learning the material
\item[-] Rely on AI for exam preparation without independent practice
\end{itemize}

\textbf{Remember:} The goal is to develop YOUR ability to work with probability distributions and understand statistical inference. AI is a tool to enhance learning, not replace it.
\end{tcolorbox}

\section*{Study Tips}

\begin{enumerate}
\item \textbf{Practice with R:} Work through all the lab exercises before attempting this assignment
\item \textbf{Verify by hand:} For simple problems, calculate by hand first, then check with R
\item \textbf{Understand, don't memorize:} Focus on understanding when to use each distribution
\item \textbf{Draw pictures:} Sketch normal curves and shade regions for probability problems
\item \textbf{Check your answers:} Probabilities must be between 0 and 1; CIs should make sense in context
\end{enumerate}

\section*{Resources}

\begin{itemize}
\item \textbf{Textbook:} Pagano \& Gauvreau, Chapters 7 and 8
\item \textbf{Lecture Slides:} Week 3 slides with examples
\item \textbf{Lab Materials:} Week 3 lab with step-by-step R code
\item \textbf{R Documentation:} \texttt{?dbinom}, \texttt{?pnorm}, \texttt{?t.test}
\item \textbf{Office Hours:} Wednesday 10:00-11:30am or by appointment
\end{itemize}

\end{document}
